% import custom functions
% Use US Letter paper, change to a4paper for A4
\documentclass[letterpaper]{article}

\usepackage{hyperref} % Required for links
\usepackage{ragged2e}
\usepackage{tikz}
\usetikzlibrary{positioning,backgrounds,shadows,snakes}

\usepackage[absolute,overlay]{textpos}
\setlength{\TPHorizModule}{1cm} % Left margin
\setlength{\TPVertModule}{1cm} % Top margin
\setlength{\parindent}{0pt} % Disable paragraph indentation

\usepackage{color}
\usepackage{amssymb}
\usepackage{calc}



%-----------------------------------------------------------------------------
%	MISC
%-----------------------------------------------------------------------------
\newcommand{\sidebarwidth}{6}

% this leaves list unchanged and returns \head and \tail
% there head is the first element and tail is all the rest
\def\spliceList#1{\expandafter\spliceListAux #1,\END}

\def\spliceListAux#1,#2\END{%
	\def\head{#1}
	\ifx\END#2\END
		\def\tail{}% tail is empty
	\else
		% remove trailing comma
		\expandafter\spliceListAuxAux #2\END
	\fi
}
\def\spliceListAuxAux#1,\END{\def\tail{#1}}



%-----------------------------------------------------------------------------
%	COLORS
%-----------------------------------------------------------------------------
\definecolor{leftsidecolor}  {HTML}{BDBDBD}
\definecolor{rightsidecolor} {HTML}{FFFFFF}
\definecolor{accentcolor}    {HTML}{607D8B}
\definecolor{darkcolor}      {HTML}{616161}



%------------------------------------------------------------------------------
%	MAIN PAGE LAYOUT
%------------------------------------------------------------------------------
\newcommand{\topoffset}{2}
\newcommand{\settopoffset}[1]{\renewcommand{\topoffset}{#1} }

\newcommand{\job}{JobTitle}
\newcommand{\setjob}[1]{\renewcommand{\job}{#1} }

\newcommand{\name}{YourName}
\newcommand{\setname}[2]{\renewcommand{\name}{{\sc #1 #2}} }

\def\slist{2,3,4,5}

\newcommand{\namesection}{
	\vspace{0.5cm}
	{\Huge \name} \\
	{\Large\job}
	\vspace{1cm}
}
\newcommand{\drawframe}{
	% Set the color of the left sidebar
	\begin{textblock}{1}(0, \topoffset)
		\begin{tikzpicture}[background rectangle/.style={fill=leftsidecolor},
			show background rectangle]
			\fill [leftsidecolor] (0,0) rectangle (8, 30){};
		\end{tikzpicture}
	\end{textblock}

	% Set the color of the right sidebar
	\begin{textblock}{1}(7.2, \topoffset)
		\begin{tikzpicture}[background rectangle/.style={fill=rightsidecolor},
			show background rectangle]
			\fill [rightsidecolor] (0,0) rectangle (20, 30){};
		\end{tikzpicture}
	\end{textblock}
}



%------------------------------------------------------------------------------
%	TITLES
%------------------------------------------------------------------------------
\newcommand{\jobtitle}[2]{
	\centerline{\textbf{\large{#1 - #2}}}
	\vspace{4pt}
}
\newcommand{\customsectionleft}[1]{
	\vspace{8pt}
	\centerline{\Large \textbf{\uppercase{#1}}}
	\rule{\textwidth}{1pt}
}
\newcommand{\customsectionright}[1]{
	\vspace{8pt}
	\centerline{\Large \textbf{\uppercase{#1}}}
	\rule{\textwidth}{1pt}
	\vspace{8pt}
}



%------------------------------------------------------------------------------
%	INFO
%------------------------------------------------------------------------------
\newcommand{\baseinfodir}{icons/}
\newcommand{\info}[2]{
	\begin{tikzpicture}
		\node (img) {\includegraphics[width=15pt] {\baseinfodir{#1}.png}};
		\node [right = 0.1cm of img] {#2};
	\end{tikzpicture}
}



%------------------------------------------------------------------------------
%	SKILLS
%------------------------------------------------------------------------------
\newcommand{\skillheight}{0.5}
\newcommand{\skilloffset}{0.15}
\newcommand{\skill}[2]{%
	{#1} \\
	\begin{tikzpicture}
		\fill [accentcolor] (0,\skilloffset) rectangle (#2,\skillheight);
		\draw [darkcolor]   (0,\skilloffset) rectangle (\sidebarwidth,\skillheight);
	\end{tikzpicture}
}

\def\myscale{0.35}
\newcommand{\showdots}[2][10]{%
	\begin{tikzpicture}
		\foreach \x in {1,...,#1} {%
			\ifnum \x > #2
				\draw[fill=leftsidecolor] (\x*\myscale,-1) circle (0.1cm);
			\else
				\draw[fill=black] (\x*\myscale,-1) circle (0.1cm);
			\fi
		}
	\end{tikzpicture}%
}

\newcommand{\drawslist}{%
	\spliceList{\slist}
	\head{}
	\foreach \s in \tail {\textbullet{}{ }\s{ }}
}
% vim: set cc=80


\settopoffset{0}
\setname{Liam}{Beguin}
\setjob{D\'eveloppeur logiciel junior}
\def\slist{Linux, Git, Syst\`emes embarqu\'es, Yocto, R\'eseaux de neurones, C, Bash,
		   \LaTeX, Kernel device drivers, Syst\`emes a\'eriens autonomes, Robotique,
		   Altium Designer}

% Set pdf metadata
\hypersetup{
	pdftitle    = {\name: \job},
	pdfsubject  = {\job},
	pdfauthor   = {\name},
	pdfcreator  = {\name},
	pdfkeywords = {\slist},
}

\begin{document}
	\thispagestyle{empty} % remove page number
	\drawframe % set background colors

	%--------------------------------------------------------------------------
	%	LEFT COLUMN
	%--------------------------------------------------------------------------

	\begin{textblock}{\sidebarwidth}(0.5, 0.5)
		\namesection

		\customsectionleft{Informations}

		\info{location1}{Montr\'eal, CA}
		\\ \info{social}{@liambeguin}
		\\ \info{language}{fran\c cais, anglais}
		\\ \info{phone1}{(514) 430-8113}
		\\ \info{email1}{\href{mailto:liambeguin@gmail.com}{liambeguin@gmail.com}}
		% \\ \info{car}{driver's license, no car}
		\vspace{0.6cm}

		\customsectionleft{Sommaire}

		Au cours de mon travail en tant que d\'eveloppeur logiciel, j'ai pu
		perfectionner mes connaissances en programmation (C, Bash, Python) et
		apprendre beaucoup sur les concepts d'architecture logicielles embarqu\'ee.
		Prochaimenemt, j'aimerai am\'eliorer mes connaissances sur le noyaux Linux
		et commencer a travailler avec la comunaut\'e. Au cours de mes exp\'eriences,
		je me suis aperçu que je pouvais m'adapter rapidement \`a de nouveaux
		objectifs et contraintes fix\'es par l'industrie de pointe. Je suis
		reconnu pour ma curiosité scientifique et mon travail autonome.
		\vspace{1cm}

		\customsectionleft{Comp\'etences}

		\begin{tabular}{ll}
			Bash          & \showdots{7} \\
			Yocto         & \showdots{6} \\
			Git           & \showdots{5} \\
			C             & \showdots{7} \\
			Python        & \showdots{6} \\
			Linux kernel  & \showdots{3} \\
			Vim           & \showdots{7} \\
		\end{tabular} \\ \\

		\drawslist{}

	\end{textblock}



	%--------------------------------------------------------------------------
	%	RIGHT COLUMN
	%--------------------------------------------------------------------------
	\begin{textblock}{13.5}(7.5, 0.5)
		\customsectionright{Exp\'eriences}

		\jobtitle{Xiphos Systems Corporation}{Montr\'eal, CA}
		D\'edi\'e essentiellement au domaine spatial, {\em Xiphos systems}
		offre des solutions hybride logique et logicielle, fournissant par exemple
		des capacit\'es de contr\^ole pour diverses charges utiles.

		\begin{description}
			\item[Mots cl\'es:] Yocto, Linux shell, Linux kernel, U-Boot, Git, Xilinx zynq, gcc
			\item[D\'eveloppeur logiciel junior - Mi-temps] \dotfill 05/2015 \`a \emph{aujourd'hui}
				\subitem{D\'eveloppement de logiciel embarqu\'e}
				\subitem{Maintenance de distribution linux}
			\item[D\'eveloppeur logiciel junior - Stage] \dotfill 09/2014 \`a 04/2015
				\subitem{Assurance qualit\'e}
		\end{description}

		\jobtitle{HiKoB}{Grenoble, FR}
		La mission de {\em HiKoB} est de d\'evelopper des syst\`emes d'acquisition
		sans fil, multipoints et autonomes dans le but de capturer l'information
		directement sur le terrain.

		\begin{description}
			\item[Mots cl\'es:] Altium designer, STM32, Git, openlab, C, IMU
			\item[d\'eveloppeur logiciel junior - Stage] \dotfill 04/2013 \`a 07/2013
				\subitem{Conception de cartes \'electroniques (PCB)}
				\subitem{D\'eveloppement de logiciel embarqu\'e}
		\end{description}

		\customsectionright{Projets}

		\jobtitle{Dronolab}{Montr\'eal, CA}
		L'objectif premier de ce club \'etudiant est de d\'evelopper et op\'erer
		des drones autonomes. Ces multirotors peuvent ensuite etre utilises dans
		des missions de reconnaissance a\'eriennes dans des applications civiles.

		\begin{description}
			\item[Mots cl\'es:] Altium designer, STM32, Control theory, Yocto, Linux, iMX6
			\item[Responsabiliti\'es] \hfill
				\subitem{Conception de cartes \'electroniques (PCB)}
				\subitem{D\'eveloppement du code de l'autopilote et de la charge utile}
				\subitem{Membre de l'\'equipe de competition \`a \emph{Unmanned Systems Canada} (2014, 2015)}
				\subitem{Membre de l'\'equipe de competition \`a \emph{AUVSI Seafarer} (2015)}
		\end{description}

		\customsectionright{Education}

			\centerline{\textbf{\large{Baccalaur\'eat en g\'enie \'electrique}}}
			\'Ecole de Technologie Sup\'erieure (ETS) \textbullet{} Montr\'eal, CA \dotfill 09/2013 \`a 12/2017 \\

			\centerline{\textbf{\large{Dipl\^ome d'\'etudes coll\'egiales en \'electronique et en informatique}}}
			Universit\'e Joseph Fourier \textbullet{} Grenoble, FR \dotfill 09/2011 \`a 07/2013 \\
	\end{textblock}
\end{document}
% vim: cc=80
